% Part: first-order-logic
% Chapter: natural-deduction
% Section: proving-things

\documentclass[../../../include/open-logic-section]{subfiles}

\begin{document}

\iftag{FOL}
      {\olfileid{fol}{ntd}{pro}}
      {\olfileid{pl}{ntd}{pro}}

\olsection{نمونه‌هایی از \usetoken{P}{derivation}}

\begin{ex}
بیایید !!a{derivation} از !!{sentence}~$(!A \land !B) \lif !A$ به دست
دهیم.

کار را با نوشتن نتیجهٔ مطلوب در پایین !!{derivation} آغاز می‌کنیم.
\begin{prooftree}
\AxiomC{}
\UnaryInfC{$(!A\land !B) \lif !A$}
\end{prooftree}

سپس باید دریابیم چه نوع استنتاجی می‌تواند به !!a{sentence} با این صورت
بینجامد. !!{main operator} نتیجه $\lif$ است، پس می‌کوشیم با استفاده از
قاعدهٔ \Intro{\lif} به نتیجه برسیم. بهتر است همچنان که پیش می‌روید
فرض‌های دخیل را بنویسید و قواعد استنتاج را برچسب بزنید تا به‌آسانی دیده
شود آیا همهٔ فرض‌ها در پایان برهان !!{discharged} شده‌اند یا نه.
\begin{prooftree}
\AxiomC{$\Discharge{!A \land !B}{1}$}
\DeduceC{$!A$}
\DischargeRule{\Intro{\lif}}{1}
\UnaryInfC{$(!A\land !B) \lif !A$}
\end{prooftree}

اکنون باید گام‌های میان فرض $!A \land !B$ و $!A$ را پر کنیم. چون تنها
با یک رابط، یعنی $\land$، سروکار داریم، باید از قاعدهٔ حذف $\land$
استفاده کنیم. این کار برهان زیر را به دست می‌دهد:
\begin{prooftree}
\AxiomC{$\Discharge{!A \land !B}{1}$}
\RightLabel{\Elim{\land}}
\UnaryInfC{$!A$}
\DischargeRule{\Intro{\lif}}{1}
\UnaryInfC{$(!A\land !B) \lif !A$}
\end{prooftree}
اکنون !!{derivation} درستی از $(!A \land !B) \lif !A$ داریم.
\end{ex}

\begin{ex}
حال بیایید !!a{derivation} از $(\lnot !A \lor !B)
\lif (!A \lif !B)$ به دست دهیم.

کار را با نوشتن نتیجهٔ مطلوب در پایین !!{derivation} آغاز می‌کنیم.
\begin{prooftree}
\AxiomC{}
\UnaryInfC{$(\lnot !A \lor !B) \lif (!A \lif !B)$}
\end{prooftree}
برای یافتن قاعده‌ای منطقی که بتواند این نتیجه را به ما بدهد، به رابط‌های
منطقیِ نتیجه نگاه می‌کنیم: $\lnot$، $\lor$ و $\lif$. فعلاً فقط نخستین
رخداد $\lif$ برایمان مهم است، زیرا این رخداد !!{main operator}
!!{sentence} در پایین اشتقاق است، در حالی که $\lnot$، $\lor$ و
رخداد دوم $\lif$ در دامنهٔ رابط دیگری قرار دارند؛ پس بعداً به آن‌ها
می‌پردازیم. بنابراین با قاعدهٔ \Intro{\lif} آغاز می‌کنیم. کاربردی درست
باید چنین باشد:
\begin{prooftree}
\AxiomC{$\Discharge{\lnot !A \lor !B}{1}$}
\DeduceC{$!A \lif !B$}
\DischargeRule{\Intro{\lif}}{1}
\UnaryInfC{$(\lnot !A \lor !B) \lif (!A \lif !B)$}
\end{prooftree}
برای ادامه دو امکان داریم. یا می‌توانیم از پایین به بالا کار را ادامه
دهیم و به‌دنبال کاربرد دیگری از قاعدهٔ \Intro{\lif} بگردیم، یا می‌توانیم
از بالا به پایین پیش برویم و قاعدهٔ \Elim{\lor} را اعمال کنیم. راه دوم
را برمی‌گزینیم. فرض $\lnot !A \lor !B$ را به‌عنوان چپ‌ترین مقدمهٔ
\Elim{\lor} به کار می‌بریم. برای کاربرد معتبر \Elim{\lor}، دو مقدمهٔ
دیگر باید با نتیجهٔ $!A \lif !B$ یکسان باشند، اما هر یک را می‌توان
به‌ترتیب از فرضی دیگر، یعنی یکی از دو منفصلِ $\lnot !A \lor !B$، اشتقاق
کرد. پس !!{derivation} ما چنین خواهد بود:
\begin{prooftree}
\AxiomC{$\Discharge{\lnot !A \lor !B}{1}$}
\AxiomC{$\Discharge{\lnot !A}{2}$}
\DeduceC{$!A \lif !B$}
\AxiomC{$\Discharge{!B}{2}$}
\DeduceC{$!A \lif !B$}
\DischargeRule{\Elim{\lor}}{2}
\TrinaryInfC{$!A \lif !B$}
\DischargeRule{\Intro{\lif}}{1}
\UnaryInfC{$(\lnot !A \lor !B) \lif (!A \lif !B)$}
\end{prooftree}

در هر یک از دو شاخهٔ سمت راست می‌خواهیم $!A
\lif !B$ را !!{derive} کنیم، که بهترین راهش استفاده از \Intro{\lif} است.
\begin{prooftree}
\AxiomC{$\Discharge{\lnot !A \lor !B}{1}$}
\AxiomC{$\Discharge{\lnot !A}{2}, \Discharge{!A}{3}$}
\DeduceC{$!B$}
\DischargeRule{\Intro{\lif}}{3}
\UnaryInfC{$!A \lif !B$}
\AxiomC{$\Discharge{!B}{2}, \Discharge{!A}{4}$}
\DeduceC{$!B$}
\DischargeRule{\Intro{\lif}}{4}
\UnaryInfC{$!A \lif !B$}
\DischargeRule{\Elim{\lor}}{2}
\TrinaryInfC{$!A \lif !B$}
\DischargeRule{\Intro{\lif}}{1}
\UnaryInfC{$(\lnot !A \lor !B) \lif (!A \lif !B)$}
\end{prooftree}

برای دو بخش مفقود !!{derivation}، به !!{derivation}sی از $!B$ از
$\lnot !A$ و $!A$ در وسط، و از $!A$ و $!B$ در سمت چپ نیاز داریم. نخست
مورد اول را بررسی کنیم. $\lnot !A$ و $!A$ دو مقدمهٔ \Elim{\lnot} هستند:
\begin{prooftree}
\AxiomC{$\Discharge{\lnot !A}{2}$}
\AxiomC{$\Discharge{!A}{3}$}
\RightLabel{\Elim{\lnot}}
\BinaryInfC{$\lfalse$}
\DeduceC{$!B$}
\end{prooftree}
با استفاده از \FalseInt می‌توانیم $!B$ را به‌عنوان نتیجه به دست آوریم
و شاخه را کامل کنیم.
\begin{prooftree}
\AxiomC{$\Discharge{\lnot !A \lor !B}{1}$}
\AxiomC{$\Discharge{\lnot !A}{2}$}
\AxiomC{$\Discharge{!A}{3}$}
\RightLabel{\Elim{\lnot}}
\BinaryInfC{$\lfalse$}
\RightLabel{\FalseInt}
\UnaryInfC{$!B$}
\DischargeRule{\Intro{\lif}}{3}
\UnaryInfC{$!A \lif !B$}
\AxiomC{$\Discharge{!B}{2}, \Discharge{!A}{4}$}
\DeduceC{$!B$}
\DischargeRule{\Intro{\lif}}{4}
\UnaryInfC{$!A \lif !B$}
\DischargeRule{\Elim{\lor}}{2}
\TrinaryInfC{$!A \lif !B$}
\DischargeRule{\Intro{\lif}}{1}
\UnaryInfC{$(\lnot !A \lor !B) \lif (!A \lif !B)$}
\end{prooftree}

اکنون به راست‌ترین شاخه نگاه کنیم. اینجا باید توجه داشت که تعریف
!!{derivation}، \emph{رفع فرض‌ها را مجاز می‌شمارد} اما آن را
\emph{الزامی نمی‌کند}. به عبارت دیگر، اگر بتوانیم $!B$ را از یکی از
فرض‌های $!A$ و $!B$ بدون استفاده از دیگری اشتقاق کنیم، اشکالی ندارد.
و !!{derive} کردن $!B$ از~$!B$ بدیهی است: $!B$ به‌تنهایی چنین
!!a{derivation}ی است و هیچ استنتاجی لازم نیست. پس می‌توانیم فرض~$!A$ را
به‌سادگی حذف کنیم.
\begin{prooftree}
\AxiomC{$\Discharge{\lnot !A \lor !B}{1}$}
\AxiomC{$\Discharge{\lnot !A}{2}$}
\AxiomC{$\Discharge{!A}{3}$}
\RightLabel{\Elim{\lnot}}
\BinaryInfC{$\lfalse$}
\RightLabel{\FalseInt}
\UnaryInfC{$!B$}
\DischargeRule{\Intro{\lif}}{3}
\UnaryInfC{$!A \lif !B$}
\AxiomC{$\Discharge{!B}{2}$}
\RightLabel{\Intro{\lif}}
\UnaryInfC{$!A \lif !B$}
\DischargeRule{\Elim{\lor}}{2}
\TrinaryInfC{$!A \lif !B$}
\DischargeRule{\Intro{\lif}}{1}
\UnaryInfC{$(\lnot !A \lor !B) \lif (!A \lif !B)$}
\end{prooftree}
توجه کنید که در !!{derivation} تکمیل‌شده، راست‌ترین استنتاج
\Intro{\lif} در واقع هیچ فرضی را رفع نمی‌کند.
\end{ex}

\begin{ex}
تا اینجا به قاعدهٔ \FalseCl{} نیاز نداشته‌ایم. این قاعده از آن حیث
ویژه است که اجازه می‌دهد فرضی را رفع کنیم که !!{formula}ی فرعی از نتیجهٔ
قاعده نیست. این قاعده ارتباط نزدیکی با قاعدهٔ \FalseInt{} دارد. در
واقع، قاعدهٔ \FalseInt{} حالت خاصی از قاعدهٔ \FalseCl{} است---منطقی به
نام «منطق شهودی» وجود دارد که در آن فقط \FalseInt{} مجاز است. قاعدهٔ
\FalseCl{} آخرین چاره است، هنگامی که هیچ روش دیگری کار نمی‌کند. برای
نمونه، فرض کنید می‌خواهیم $!A \lor \lnot !A$ را !!{derive} کنیم. راهبرد
معمول ما تلاش برای !!{derive} کردن $!A \lor \lnot !A$ با استفاده از
$\Intro{\lor}$ است. اما این امر مستلزم آن است که یا $!A$ یا $\lnot !A$
را بدون هیچ فرضی !!{derive} کنیم، و این کار شدنی نیست. \FalseCl{} به
دادمان می‌رسد!
\begin{prooftree}
  \AxiomC{$\Discharge{\lnot(!A \lor \lnot !A)}{1}$}
  \DeduceC{$\lfalse$}
  \DischargeRule{\FalseCl}{1}
  \UnaryInfC{$!A \lor \lnot !A$}
\end{prooftree}
اکنون به‌دنبال !!a{derivation} از $\lfalse$ از $\lnot(!A \lor
\lnot !A)$ هستیم. چون $\lfalse$ نتیجهٔ $\Elim{\lnot}$ است، شاید آن را
امتحان کنیم:
\begin{prooftree}
  \AxiomC{$\Discharge{\lnot(!A \lor \lnot !A)}{1}$}
  \DeduceC{$\lnot !A$}
  \AxiomC{$\Discharge{\lnot(!A \lor \lnot !A)}{1}$}
  \DeduceC{$!A$}
  \RightLabel{\Elim{\lnot}}
  \BinaryInfC{$\lfalse$}
  \DischargeRule{\FalseCl}{1}
  \UnaryInfC{$!A \lor \lnot !A$}
\end{prooftree}
راهبردمان برای یافتن !!a{derivation} از~$\lnot !A$ کاربرد
~$\Intro{\lnot}$ را اقتضا می‌کند:
\begin{prooftree}
  \AxiomC{$\Discharge{\lnot(!A \lor \lnot !A)}{1}, \Discharge{!A}{2}$}
  \DeduceC{$\lfalse$}
  \DischargeRule{\Intro{\lnot}}{2}
  \UnaryInfC{$\lnot !A$}
  \AxiomC{$\Discharge{\lnot(!A \lor \lnot !A)}{1}$}
  \DeduceC{$!A$}
  \RightLabel{\Elim{\lnot}}
  \BinaryInfC{$\lfalse$}
  \DischargeRule{\FalseCl}{1}
  \UnaryInfC{$!A \lor \lnot !A$}
\end{prooftree}
در اینجا می‌توانیم $\lfalse$ را به‌آسانی با اعمال $\Elim{\lnot}$ بر فرض
$\lnot(!A \lor \lnot !A)$ و $!A \lor \lnot !A$ به دست آوریم؛ دومی
از فرض تازهٔ $!A$ با~$\Intro{\lor}$ نتیجه می‌شود:
\begin{prooftree}
  \AxiomC{$\Discharge{\lnot(!A \lor \lnot !A)}{1}$}
  \AxiomC{$\Discharge{!A}{2}$}
  \RightLabel{\Intro{\lor}}
  \UnaryInfC{$!A \lor \lnot !A$}
  \RightLabel{\Elim{\lnot}}
  \BinaryInfC{$\lfalse$}
  \DischargeRule{\Intro{\lnot}}{2}
  \UnaryInfC{$\lnot !A$}
  \AxiomC{$\Discharge{\lnot(!A \lor \lnot !A)}{1}$}
  \DeduceC{$!A$}
  \RightLabel{\Elim{\lnot}}
  \BinaryInfC{$\lfalse$}
  \DischargeRule{\FalseCl}{1}
  \UnaryInfC{$!A \lor \lnot !A$}
\end{prooftree}
در سمت راست از همان راهبرد استفاده می‌کنیم، جز آنکه $!A$ را
با~\FalseCl به دست می‌آوریم:
\begin{prooftree}
  \AxiomC{$\Discharge{\lnot(!A \lor \lnot !A)}{1}$}
  \AxiomC{$\Discharge{!A}{2}$}
  \RightLabel{\Intro{\lor}}
  \UnaryInfC{$!A \lor \lnot !A$}
  \RightLabel{\Elim{\lnot}}
  \BinaryInfC{$\lfalse$}
  \DischargeRule{\Intro{\lnot}}{2}
  \UnaryInfC{$\lnot !A$}
  \AxiomC{$\Discharge{\lnot(!A \lor \lnot !A)}{1}$}
  \AxiomC{$\Discharge{\lnot !A}{3}$}
  \RightLabel{\Intro{\lor}}
  \UnaryInfC{$!A \lor \lnot !A$}
  \RightLabel{\Elim{\lnot}}
  \BinaryInfC{$\lfalse$}
  \DischargeRule{\FalseCl}{3}
  \UnaryInfC{$!A$}
  \RightLabel{\Elim{\lnot}}
  \BinaryInfC{$\lfalse$}
  \DischargeRule{\FalseCl}{1}
  \UnaryInfC{$!A \lor \lnot !A$}
\end{prooftree}
\end{ex}

\begin{prob}
!!{derivation}sی به دست دهید که موارد زیر را نشان دهند:
\begin{enumerate}
\item $!A \land (!B \land !C) \Proves (!A \land !B) \land !C$.
\item $!A \lor (!B \lor !C) \Proves (!A \lor !B) \lor !C$.
\item $!A \lif (!B \lif !C) \Proves !B \lif (!A \lif !C)$.
\item $!A \Proves \lnot\lnot !A$.
\end{enumerate}
\end{prob}

\begin{prob}
!!{derivation}sی به دست دهید که موارد زیر را نشان دهند:
\begin{enumerate}
\item $(!A \lor !B) \lif !C \Proves !A \lif !C$.
\item $(!A \lif !C) \land (!B \lif !C) \Proves (!A \lor !B) \lif !C$.
\item $\Proves \lnot(!A \land \lnot !A)$.
\item $!B \lif !A \Proves \lnot !A \lif \lnot !B$.
\item $\Proves (!A \lif \lnot !A) \lif \lnot !A$.
\item $\Proves \lnot(!A \lif !B) \lif \lnot !B$.
\item $!A \lif !C \Proves \lnot (!A \land \lnot !C)$.
\item $!A \land \lnot !C \Proves \lnot (!A \lif !C)$.
\item $!A \lor !B, \lnot !B \Proves !A$.
\item $\lnot !A \lor \lnot !B \Proves \lnot(!A \land !B)$.
\item $\Proves (\lnot !A \land \lnot !B) \lif\lnot(!A \lor !B)$.
\item $\Proves \lnot(!A \lor !B) \lif (\lnot !A \land \lnot !B)$.
\end{enumerate}
\end{prob}

\begin{prob}
!!{derivation}sی به دست دهید که موارد زیر را نشان دهند:
\begin{enumerate}
\item $\lnot(!A \lif !B) \Proves !A$.
\item $\lnot(!A \land !B) \Proves \lnot !A \lor \lnot !B$.
\item $!A \lif !B \Proves \lnot !A \lor !B$.
\item $\Proves \lnot \lnot !A \lif !A$.
\item $!A \lif !B, \lnot !A \lif !B \Proves !B$.
\item $(!A \land !B) \lif !C \Proves (!A \lif !C) \lor (!B \lif !C)$.
\item $(!A \lif !B) \lif !A \Proves !A$.
\item $\Proves (!A \lif !B) \lor (!B \lif !C)$.
\end{enumerate}
(همهٔ این‌ها به قاعدهٔ~$\FalseCl$ نیاز دارند.)
\end{prob}

\end{document}
